% =====================================================================================
% VERSION UNE PAGE de la note de themes, pour Hugo Rapior avant le point RH du 4 septembre.
% La version longue (2026-09-04_themes_travail.tex) garde les motifs detailles, le bloc sur
% ce que les themes ouvrent pour le cabinet, et l'emplacement des missions Nexialog.
%
% CE QUE CETTE VERSION DOIT FAIRE : se lire en deux minutes et tenir sur UNE page. Trois
% rubriques, dans l'ordre demande : developper, porter, staffable. Rien d'autre.
%
% DEUX CONTRAINTES DE FORME, VOLONTAIRES :
%   - aucun tiret cadratin. Le babel francais en met un en puce d'itemize par defaut, d'ou
%     le label redefini plus bas. A verifier sur le PDF, jamais sur la source ;
%   - aucun chiffre qui ne se releve sur le depot du memoire.
% =====================================================================================
\documentclass[11pt,a4paper]{article}

\usepackage[utf8]{inputenc}
\usepackage[T1]{fontenc}
\usepackage{lmodern}
\usepackage[french]{babel}
\usepackage[top=1.6cm,bottom=1.5cm,left=1.9cm,right=1.9cm]{geometry}
\usepackage{array}
\usepackage{enumitem}
\usepackage{xcolor}
\usepackage{titlesec}

% Charte Nexialog : memes roles que style_nexialog.py, aucun hexadecimal ailleurs qu'ici.
\definecolor{encre}{HTML}{1B1E30}
\definecolor{accent}{HTML}{A6002E}
\definecolor{bleu}{HTML}{204993}
\definecolor{grisclair}{HTML}{F2F2F2}

\color{encre}
\titleformat{\section}{\normalfont\bfseries\color{bleu}}{}{0pt}{}
\titlespacing{\section}{0pt}{10pt}{3pt}
% Puce explicite : le babel francais mettrait un tiret cadratin.
\setlist[itemize]{label=\textcolor{accent}{\textbullet},itemsep=2pt,topsep=2pt,
                  leftmargin=1.1em,parsep=0pt}

\pagestyle{empty}

\begin{document}
\small

{\Large\bfseries\color{bleu} Thèmes de travail souhaités}\\[3pt]
{\color{accent}\rule{\textwidth}{1.1pt}}\\[3pt]
{\footnotesize Kélian Kaddouri \hfill Pour Hugo Rapior, avant le point RH du 4 septembre
2026}

\vspace{5pt}

\noindent\textbf{En une phrase.} Devenir le référent quantitatif du cabinet sur le risque
cyber, en restant capable de tenir les sujets non-vie classiques, provisionnement et capital,
sans temps de montée en charge. C'est chez Nexialog que je veux construire cette trajectoire.

\section{Ce que je veux développer}

\begin{itemize}
\item \textbf{R\&D cyber, et nommément le Lab Cyber.} Sujet jeune, donnée rare, aucune
      méthode de référence établie, et une obligation réglementaire depuis janvier 2025 sans
      qu'aucun module de formule standard ne la couvre : la demande existe avant la méthode.
      C'est le périmètre annoncé du Lab, quantification du risque cyber en lien avec DORA,
      NIS2 et Solvabilité II, et c'est exactement ce sur quoi j'ai travaillé un an.
      \textit{Ce que j'apporte} : un mémoire qui quantifie un besoin de capital en pilier~2
      au titre de DORA par une cascade dirigée entre les cinq piliers, fréquence et sévérité,
      Monte-Carlo, calibration valeurs extrêmes, identification partielle de la matrice de
      propagation, et chaque nombre publié rattaché au script qui le calcule.
\item \textbf{Non-vie.} La branche où vivent les outils que je maîtrise, fréquence et
      sévérité, queues lourdes, charge annuelle agrégée. Le cyber en est une composante, pas
      une exception, et c'est ce qui rend le profil plaçable en dehors de sa spécialité.
      \textit{Ce que j'apporte} : fréquence et sévérité modélisées séparément, tests
      d'adéquation, surdispersion, mesures de risque sur charge agrégée et allocation par
      Euler.
\item \textbf{Réassurance et accumulation.} Le prolongement direct de mon sujet : dès qu'on
      parle accumulation et cause commune, la question devient une question de structure de
      traité, par risque ou par événement, et la cession est le premier levier qu'un client
      actionne avant le capital. Y compris lorsque cela se traite comme un programme de
      cession à l'intérieur d'un dossier non-vie, sans mandat de réassurance dédié.
      \textit{Ce que j'apporte} : accumulation par cause commune, dépendance de queue,
      principe du grand saut unique, et le pont entre quantile d'un sinistre et quantile de
      la charge annuelle.
\item \textbf{Risque climatique.} C'est l'un des trois labs du pôle R\&D et l'un des thèmes
      évoqués en séance, donc un axe du cabinet et non une idée de ma part. Je ne revendique pas
      une expertise climat : je revendique une méthode qui traverse, la cascade et l'accumulation
      par cause commune étant les mêmes objets d'un péril à l'autre.
      \textit{Ce que j'apporte} : l'audit d'un préprint de cascade climatique multi-périls, dont
      j'ai mesuré les cinq notions sur mon propre modèle, ce qui a montré que l'ordre de leurs
      résultats dépend de l'indice de queue et ne se transporte donc pas tel quel. Plus
      l'ablation en échelle et le protocole de comparaison des structures de dépendance.
\end{itemize}

\section{Ce que je veux porter}

\noindent Sans en faire un axe, et sans réserve sur l'engagement.

\begin{itemize}
\item \textbf{Data science appliquée à la fraude.} Données déséquilibrées, arbitrage entre
      pouvoir prédictif et explicabilité. Le rapprochement avec le cyber est technique :
      événement rare, donnée censurée par la détection, et ce qui n'est pas vu compte autant
      que ce qui est vu.
\item \textbf{Provisionnement.} Technique et normé, et c'est le sujet qui donne accès aux
      directions techniques.
\item \textbf{Modélisation du capital.} Formule standard, SCR par module, et le pilier~2 des
      deux côtés, ORSA en assurance comme ICAAP en banque. C'est déjà le terrain de mon
      mémoire.
\item \textbf{Gestion des risques.} Cartographie, criticité, appétit au risque, articulation
      avec l'ORSA. C'est d'où vient mon sujet : la cascade entre les cinq piliers a été une
      cartographie qualitative avant d'être un modèle stochastique.
\item \textbf{IA générative.} Je l'utilise au quotidien pour produire et pour relire, et
      l'angle qui m'intéresse est de l'encadrer comme un risque de modèle, avec la même
      exigence de traçabilité que pour un modèle de capital. Je ne la revendique pas comme
      une expertise, je la porte volontiers.
\end{itemize}

\section{Staffable sans montée en charge}

\begin{itemize}
\item \textbf{Solvabilité II} : SCR par module, besoin de capital en ORSA, allocation par
      entité, et la calibration des lois de perte qui va avec, sur données incomplètes.
\item \textbf{Industrialisation en Python} : une chaîne de plus de quatre-vingt-dix scripts
      versionnés, modules partagés, sorties reproductibles à l'octet entre deux postes.
\end{itemize}

% =====================================================================================
% BLOC A REMPLIR, ET C'EST CE QUI PESERA LE PLUS. Hugo a demande en seance que les
% descriptions disent le CONTEXTE et le GAIN (exemple cite : le gain sur l'Ains Index).
% Modele par ligne, dans cet ordre : secteur ou client (anonymise si besoin) ; ce qui etait
% demande ; ce que J'AI fait ; le resultat CHIFFRE ou la decision permise.
%
% POUR TENIR SUR UNE PAGE en le decommentant, retirer l'encadre « Note en seance »
% ci-dessus : son argument est deja porte par la version deux pages.
%
% \section{Ce que j'ai livré chez Nexialog}
% \begin{itemize}
% \item \textbf{[Secteur, type de mission].} [Demande.] [Ce que j'ai fait.] [Resultat chiffre.]
% \item \textbf{[Secteur, type de mission].} [...]
% \end{itemize}
% =====================================================================================

\vspace{6pt}
\noindent{\color{accent}\rule{\textwidth}{0.6pt}}\\[1pt]
{\footnotesize\itshape Liste à valider avant le point RH du 4 septembre, pour transmission à
Axelle et Ali. Version détaillée disponible : motifs, apport pour le cabinet, et missions.}

\end{document}
